% !TeX root = ../main.tex
\section{Proofs for Section~\ref{sec:preliminaries}}
\label{app:preliminaries}


\subsection{Proof of the Quantitative Primal-Function Oracle Bounds}
\label{app:corrected-estimators}

\begin{proof}[Proof of Lemma~\ref{lem:corrected-errors}]

\paragraph{Corrected value.}
Fix $x$, and write
\[
    C:=\nabla^2_{yy}f(x,y),
    \qquad
    q:=\nabla_yf(x,y),
    \qquad
    \delta:=y-y^*(x).
\]
Taylor expansion of the inner gradient gives
\begin{equation}
    v:=q-C\delta,
    \qquad
    \|v\|
    \le
    \frac{\rho}{2}\|\delta\|^2.
    \label{eq:app-value-v-quadratic}
\end{equation}
Define
\[
    \overline C
    :=
    \int_0^1
    \nabla^2_{yy}f(x,y-t\delta)\,dt.
\]
The fundamental theorem of calculus gives
$q=\overline C\delta$, and hence
\[
    v=(\overline C-C)\delta.
\]
Both $C$ and $\overline C$ satisfy
\[
    -\ell I
    \preceq
    C,\overline C
    \preceq
    -\mu I,
\]
and therefore
\begin{equation}
    \|v\|
    \le
    (\ell-\mu)\|\delta\|.
    \label{eq:app-value-v-linear}
\end{equation}

Taylor expansion of the payoff around $y$ gives
\begin{equation}
    P(x)
    =
    f(x,y)
    -
    q^\top\delta
    +
    \frac12\delta^\top C\delta
    +
    R,
    \qquad
    |R|
    \le
    \frac{\rho}{6}\|\delta\|^3.
    \label{eq:app-value-taylor}
\end{equation}
Using $q=C\delta+v$ and completing the square yields
\begin{equation}
    P(x)-\widehat P(x,y)
    =
    R+\frac12v^\top C^{-1}v.
    \label{eq:app-value-square}
\end{equation}
Strong concavity gives
$\|C^{-1}\|\le1/\mu$.
Combining
\eqref{eq:app-value-v-quadratic} and
\eqref{eq:app-value-v-linear},
\[
    \|v\|^2
    \le
    \frac{\rho}{2}\|\delta\|^2
    (\ell-\mu)\|\delta\|.
\]
Hence
\[
    \frac12
    |v^\top C^{-1}v|
    \le
    \frac{(\kappa-1)\rho}{4}\|\delta\|^3.
\]
Together with \eqref{eq:app-value-taylor}, this gives
\[
    |\widehat P(x,y)-P(x)|
    \le
    \frac{(3\kappa-1)\rho}{12}\|\delta\|^3.
\]
The inequality
\[
    \frac{(3\kappa-1)\rho}{12}
    \le
    \frac{\kappa\rho}{4}
\]
follows from $\kappa\ge1$.

If $\|\delta\|\le\mu/\rho$, then
\eqref{eq:app-value-v-quadratic} gives
\[
    \frac12|v^\top C^{-1}v|
    \le
    \frac{\rho^2}{8\mu}\|\delta\|^4
    \le
    \frac{\rho}{8}\|\delta\|^3.
\]
Combining this with
$|R|\le\rho\|\delta\|^3/6$ proves
\[
    |\widehat P(x,y)-P(x)|
    \le
    \frac{7\rho}{24}\|\delta\|^3.
\]


\paragraph{Corrected gradient.}
Let
\[
    y_t
    :=
    y^*(x)+t\bigl(y-y^*(x)\bigr),
    \qquad
    t\in[0,1].
\]
Since $f$ is twice continuously differentiable, the fundamental
theorem of calculus gives
\begin{align}
    \nabla_xf(x,y)-\nabla_xf(x,y^*(x))
    &=
    \int_0^1
    \nabla^2_{xy}f(x,y_t)
    \bigl(y-y^*(x)\bigr)\,dt,
    \label{eq:app-gradient-x-path}
    \\
    \nabla_yf(x,y)
    &=
    \int_0^1
    \nabla^2_{yy}f(x,y_t)
    \bigl(y-y^*(x)\bigr)\,dt,
    \label{eq:app-gradient-y-path}
\end{align}
where the second identity uses
$\nabla_yf(x,y^*(x))=0$.

Substituting
\eqref{eq:app-gradient-x-path}--\eqref{eq:app-gradient-y-path}
into \eqref{eq:corrected-gradient} and using
\eqref{eq:primal-gradient} gives
\begin{align}
    \widehat g(x,y)-\nabla P(x)
    &=
    \int_0^1
    \Big[
        \nabla^2_{xy}f(x,y_t)
        -
        \nabla^2_{xy}f(x,y)
        \notag\\
    &\qquad
        +
        \nabla^2_{xy}f(x,y)
        \bigl(\nabla^2_{yy}f(x,y)\bigr)^{-1}
        \bigl(
            \nabla^2_{yy}f(x,y)
            -
            \nabla^2_{yy}f(x,y_t)
        \bigr)
    \Big]
    \notag\\
    &\qquad\qquad\times
    \bigl(y-y^*(x)\bigr)\,dt.
    \label{eq:app-gradient-cancellation}
\end{align}
This identity displays the cancellation of the linear inner error.

Hessian Lipschitz continuity and
\[
    \left\|
    \nabla^2_{xy}f(x,y)
    \bigl(\nabla^2_{yy}f(x,y)\bigr)^{-1}
    \right\|
    \le
    \kappa
\]
give
\[
    \|\widehat g(x,y)-\nabla P(x)\|
    \le
    (1+\kappa)\rho
    \int_0^1(1-t)\,dt\,
    \|y-y^*(x)\|^2.
\]
Therefore
\[
    \|\widehat g(x,y)-\nabla P(x)\|
    \le
    \frac{(1+\kappa)\rho}{2}
    \|y-y^*(x)\|^2.
\]


\paragraph{Schur-complement Hessian.}
Retain
\[
    C:=\nabla^2_{yy}f(x,y)
\]
and set
\[
    B:=\nabla^2_{xy}f(x,y),
    \qquad
    B_*:=\nabla^2_{xy}f(x,y^*(x)),
    \qquad
    C_*:=\nabla^2_{yy}f(x,y^*(x)).
\]
The inverse identity gives
\begin{equation}
\begin{aligned}
    &
    \left\|
    C^{-1}-C_*^{-1}
    \right\|
    \\
    &\qquad\le
    \frac{\rho}{\mu^2}
    \|y-y^*(x)\|.
\end{aligned}
\label{eq:app-inverse-difference}
\end{equation}

The Schur-complement products satisfy
\begin{align}
    BC^{-1}B^\top-B_*C_*^{-1}B_*^\top
    &=
    (B-B_*)C^{-1}B^\top
    \notag\\
    &\quad+
    B_*(C^{-1}-C_*^{-1})B^\top
    \notag\\
    &\quad+
    B_*C_*^{-1}(B-B_*)^\top.
    \label{eq:app-schur-telescope}
\end{align}
Hessian Lipschitz continuity,
$\|B\|,\|B_*\|\le\ell$, and
$\|C^{-1}\|,\|C_*^{-1}\|\le1/\mu$
bound the three terms on the right-hand side by
\[
    \kappa\rho\|y-y^*(x)\|,
    \qquad
    \kappa^2\rho\|y-y^*(x)\|,
    \qquad
    \kappa\rho\|y-y^*(x)\|,
\]
respectively.
Adding the $xx$-block error yields
\begin{align*}
    \|\widehat H(x,y)-\nabla^2P(x)\|
    &\le
    \rho\|y-y^*(x)\|
    +
    2\kappa\rho\|y-y^*(x)\|
    \\
    &\qquad+
    \kappa^2\rho\|y-y^*(x)\|
    \\
    &=
    (1+\kappa)^2\rho
    \|y-y^*(x)\|.
\end{align*}
This proves the Hessian estimate and completes the proof.
\end{proof}


\subsection{Proof of the Common Inexact Trust-Region Estimate}
\label{app:common-tr-step}

\begin{proof}[Proof of Lemma~\ref{lem:common-inexact-tr-step}]
Apply Lemma~\ref{lem:tr-global-optimality} with
\[
    B=H+aI.
\]
Its stationarity and positive-semidefiniteness conditions give
\[
    q(d)
    =
    -\frac12d^\top(H+aI)d-\lambda r^2
    \le
    -\frac{\lambda}{2}r^2,
\]
which proves \eqref{eq:common-model-bound}.

The Taylor bound
\eqref{eq:function-taylor} and the current endpoint errors imply
\begin{align*}
    F(x)-F(x+d)
    &\ge
    -g^\top d
    -\frac12d^\top Hd
    -\delta_0r
    -\frac{\xi}{2}r^2
    -\frac{M}{6}r^3.
\end{align*}
Since
\[
    -g^\top d-\frac12d^\top Hd
    =
    -q(d)+\frac{a}{2}r^2,
\]
the model bound gives
\[
    F(x)-F(x+d)
    \ge
    \frac{a-\xi}{2}r^2
    +
    \frac{\lambda}{2}r^2
    -
    \frac{M}{6}r^3
    -
    \delta_0r,
\]
which proves \eqref{eq:common-function-bound}.

For the candidate gradient, decompose
\begin{align*}
    g^+
    &=
    [g^+-\nabla F(x+d)]
    \\
    &\quad+
    [\nabla F(x+d)-\nabla F(x)-\nabla^2F(x)d]
    \\
    &\quad+
    [\nabla F(x)-g]
    +
    [\nabla^2F(x)-H]d
    +
    [g+Hd].
\end{align*}
The stationarity condition in
\eqref{eq:tr-global-optimality}, again with $B=H+aI$, gives
\[
    g+Hd=-(a+\lambda)d.
\]
Taking norms and using
\eqref{eq:gradient-taylor},
\eqref{eq:common-current-errors}, and
\eqref{eq:common-candidate-error} yields
\[
    \|g^+\|
    \le
    (a+\xi+\lambda)r
    +
    \frac{M}{2}r^2
    +
    \delta_0+\delta_+,
\]
which proves \eqref{eq:common-gradient-bound}.
\end{proof}

